% Generated draft: requires mathematical and source review.
\documentclass[11pt]{article}
\usepackage[margin=25mm]{geometry}
\usepackage{fontspec}
\setmainfont{Latin Modern Roman}
\usepackage{amsmath,amssymb,mathtools,unicode-math}
\setmathfont{Latin Modern Math}
\usepackage{enumitem,needspace,xurl,hyperref}
\setlength{\parindent}{0pt}
\setlength{\parskip}{5pt}
\setlength{\emergencystretch}{3em}
\newcommand{\N}{\mathbb{N}}
\newcommand{\Z}{\mathbb{Z}}
\newcommand{\Q}{\mathbb{Q}}
\newcommand{\R}{\mathbb{R}}
\newcommand{\C}{\mathbb{C}}
\newcommand{\F}{\mathbb{F}}
\newcommand{\A}{\mathbb{A}}
\newcommand{\PP}{\mathbb{P}}
\newcommand{\E}{\mathbb{E}}
\newcommand{\eps}{\varepsilon}
\newcommand{\dd}{\,\mathrm{d}}
\begin{document}
\section*{UnsolvedMath ID: 30002960}
\textbf{Source problem code:} catalogue entry
\subsection{Definitions and mathematical statement}
Use the following provisional interpretation of the supplied question: a triangulation is a finite simplicial complex whose geometric realization is homeomorphic to the sphere
\[
S^3=\{x\in\R^4:\|x\|=1\}.
\]
Let \(\mathcal T(S^3)\) denote the class of these triangulations. The supplied record does not explicitly specify finiteness or whether more general triangulations with face identifications are admitted; the intended class requires confirmation.

For \(T\in\mathcal T(S^3)\), write \(V(T)\) for its vertex set and define
\[
\chi_3(T)=\min\Bigl\{k\in\Z_{\geq 1}:\exists c:V(T)\to\{1,\ldots,k\}\ \text{such that}\ \forall\sigma\in T,\ \dim\sigma=3\Longrightarrow |c(V(\sigma))|\geq 2\Bigr\}.
\]
Thus the coloring is a vertex coloring: every tetrahedron must contain vertices of at least two colors, while monochromatic triangles are permitted. No proper-coloring requirement is imposed on edges. Define the uniform coloring number by
\[
\chi_3(S^3)=\sup_{T\in\mathcal T(S^3)}\chi_3(T)\in\Z_{\geq1}\cup\{\infty\}.
\]
The vertex-coloring convention and this supremum definition make explicit the natural interpretation of the notation in the record, but are not explicitly defined there.

\textbf{Problem.} Determine \(\chi_3(S^3)\). In particular, does there exist an integer \(k\geq1\) such that every \(T\in\mathcal T(S^3)\) has a vertex coloring with at most \(k\) colors and no monochromatic tetrahedron? If such a uniform bound exists, determine the least possible bound or establish an upper bound. The bound is required to be independent of the triangulation; each individual finite triangulation already admits a coloring using a distinct color at every vertex.

The supplied record labels this problem open. Its dated literature assessment instead concerns knot complexity and the treewidth of minimal triangulations of knot complements, so that assessment does not substantiate the status of this coloring question. The correspondence between the question and the cited report requires human verification.
\subsection{Short English statement}
For every finite simplicial triangulation of the \(3\)-sphere, color its vertices so that no tetrahedron is monochromatic; monochromatic triangles are allowed. Determine whether one fixed finite number of colors works for every triangulation, and determine the optimal number or an upper bound. This formulation uses a provisional interpretation of the triangulation class and coloring-number notation. The supplied open-status label has not been independently verified.
\subsection{Sources}
\begin{itemize}
\item Supplied local UnsolvedMath registry record, permanent ID \(30002960\), title \emph{Three-Dimensional Coloring Number of the Sphere}, in the collection \emph{Oberwolfach Reports Open Problems}; exact source problem code \texttt{catalogue entry}. The record gives proposed year \(2015\) and identifies no proposer.
\item Supplied source citation: \emph{Computational Geometric and Algebraic Topology (2016)}, Oberwolfach Reports, pp.\ 2637--2699, DOI \texttt{10.4171/owr/2015/45}. Supplied source URL: \url{https://doi.org/10.4171/owr/2015/45}. The background describes Oberwolfach Reports as documenting mathematical workshops and their open-problem discussions.
\item Supplied literature source: \emph{Oberwolfach Report 45, 2015}, \url{https://doi.org/10.4171/owr/2015/45}. Literature checked at: \texttt{2026-08-21}. This date and assessment are supplied dated curation, not independently verified. The supplied assessment labels the problem open but discusses a proposed relation between knot complexity and the treewidth of minimal triangulations of knot complements, described as a necessary condition for an algorithmic strategy, with no general diverging function \(f\) established. That assessment concerns a different mathematical target from the coloring question. The background itself requests independent verification of time-sensitive status and cited claims.
\item The registry supplies the statement-status label \texttt{corrected\_verified}, with a note that mathematical content was preserved while notation, typography, or omitted definitions were repaired for self-contained publication. This is supplied metadata, not independent verification of the present formulation or its provenance.
\end{itemize}
\end{document}
